\section{Applying Social Choice to RAG Pipelines}
\label{sec:social_choice}

Social choice theory as a framework is inherently context-dependent. The values that guide a medical decision such as the kidney allocation problem \citep{freedman2020akidneyallocation} are likely to be different from the values that determine whether a self-driving car should crash into a wall to kill its passengers or swerve to kill a pedestrian \citep{noothigattu2018voting}. Additionally, in many real-world scenarios, social choice preferences vary across cultural, demographic, and socioeconomic lines. For example, the moral preferences of people when asked the autonomous vehicle dilemma vary significantly across countries, and are strongly influenced by factors such as culture (e.g. collectivist vs individualist) and demographics (e.g. gender) \citep{awad2018moral}. 

\citet{pmlr-v235-conitzer24a} suggest that in these cases, where preferences vary strongly across groups (e.g. political topics, social norms) social choice principles should be applied with caution and consideration for the specific context and the values of the communities involved. 
